% ============================================================
% BOXES, CHAPTER/SUBSECTION BARS, AND HELPERS
% ============================================================

\newcommand{\BreakIfNotEnoughSpace}{%
  \par\penalty0\vspace{0pt plus \ZSFbreakThreshold}\penalty-100\vspace{0pt plus -\ZSFbreakThreshold}\relax
}

% Einheitlicher Top-Spacing-Kontrakt fuer Titelbalken.
% Reihenfolge ist kritisch: \addvspace MUSS vor \BreakIfNotEnoughSpace stehen,
% damit es per Max-Semantik mit dem after-skip der Vorbox kollabiert. Stuende
% es danach, faende \addvspace nur noch die Stretch-Reste (lastskip natural=0)
% und wuerde additiv arbeiten -> Vorbox-After-Skip + Bar-Before-Skip stapeln
% sich, was zu sichtbar variablem Abstand ueber Titelbalken fuehrt.
% Tcolorbox-Definitionen, die diesen Helper aufrufen, MUESSEN before skip=0pt
% setzen, damit das Spacing nicht doppelt eingerechnet wird.
\newcommand{\ZSFTitleBarTopSpacing}[1]{%
  \par
  \addvspace{#1}%
  \BreakIfNotEnoughSpace
}

\newtcolorbox{chapterbar}[1][] {
  enhanced jigsaw,
  breakable=false,
  colback=\ChapterBarColor,
  colframe=\ChapterBarColor,
  boxrule=0pt,
  boxsep=0pt,
  left=\ZSFspaceS,
  right=\ZSFspaceS,
  top=\ZSFspaceS,
  bottom=\ZSFspaceS,
  before skip=\ZSFspaceL,
  after skip=\ZSFchapterBarAfterSkip,
  fontupper=\ZSFfontChapter,
  #1
}
\newcommand{\ChapterBar}[1]{
  \begin{chapterbar}
  \textcolor{white}{#1}
  \end{chapterbar}
}

% ------------------------------------------------------------
% Dokument-Titel-Masthead (einmalig, vor dem ersten Kapitel)
% Kompakte Variante fuer Analysis: knapper Platz im 4-Spalten-Layout.
% ------------------------------------------------------------
\newtcolorbox{zsftitleheader}[1][]{
  enhanced jigsaw,
  breakable=false,
  colback=ChapterColor0,
  colframe=ChapterColor0,
  boxrule=0pt,
  boxsep=0pt,
  left=\ZSFspaceM,
  right=\ZSFspaceM,
  top=\ZSFspaceS,
  bottom=\ZSFspaceS,
  before skip=0pt,
  after skip=\ZSFspaceXS,
  halign=center,
  #1
}
\newcommand{\ZSFTitleHeader}[2]{%
  \penalty-500
  \begin{zsftitleheader}
    \sffamily
    {\color{white}\fontsize{11bp}{12.5bp}\bfseries\selectfont #1\par}%
    \vspace{1pt}%
    {\color{white!88}\fontsize{6.7bp}{7.5bp}\itshape\selectfont von~#2\par}%
  \end{zsftitleheader}
  \penalty10000
}

\newtcolorbox{subsectionbar}[1][] {
  enhanced jigsaw,
  breakable=false,
  colback=\SubsectionBarColor,
  colframe=\SubsectionBarColor,
  boxrule=0pt,
  boxsep=0pt,
  left=\ZSFspaceS,
  right=\ZSFspaceS,
  top=\ZSFspaceS,
  bottom=\ZSFspaceS,
  % before skip wird vom Wrapper (\ZSFTitleBarTopSpacing) verwaltet.
  before skip=0pt,
  after skip=\ZSFsubsectionAfterGap,
  fontupper=\ZSFfontSection,
  #1
}
\makeatletter
\newcommand{\SubsectionBar}{\@ifstar{\SubsectionBarStar}{\SubsectionBarNumbered}}
\newcommand{\SubsectionBarNumbered}[2][]{
  \ZSFTitleBarTopSpacing{\ZSFsubsectionBarBeforeSkip}
  \penalty-500
  \ZSFCollapseAfterChapterBar
  \refstepcounter{subsection}
  \if\relax\detokenize{#1}\relax\else\label{#1}\fi
  \begin{subsectionbar}
  \textcolor{white}{\thesubsection. #2}
  \end{subsectionbar}
  \penalty10000
  \ZSFMarkAfterSubsectionBar
}
\newcommand{\SubsectionBarStar}[2][]{
  \ZSFTitleBarTopSpacing{\ZSFsubsectionBarBeforeSkip}
  \penalty-500
  \ZSFCollapseAfterChapterBar
  \if\relax\detokenize{#1}\relax\else\label{#1}\fi
  \begin{subsectionbar}
  \textcolor{white}{#2}
  \end{subsectionbar}
  \penalty10000
  \ZSFMarkAfterSubsectionBar
}
% Variante: Abschnitt auf neuer Spalte beginnen (ersetzt rohes \columnbreak)
\newcommand{\SubsectionBarOnNewColumn}{\@ifstar{\SubsectionBarOnNewColumnStar}{\SubsectionBarOnNewColumnNumbered}}
\newcommand{\SubsectionBarOnNewColumnNumbered}[2][]{%
  \columnbreak
  \SubsectionBarNumbered[#1]{#2}%
}
\newcommand{\SubsectionBarOnNewColumnStar}[2][]{%
  \columnbreak
  \SubsectionBarStar[#1]{#2}%
}
\makeatother

% Gemeinsamer Basis-Stil fuer alle Titelleisten
\tcbset{zsftitlebar/.style={
  fonttitle=\ZSFfontBoxTitle\strut,
  halign title=center,
  title style={minimum height=\ZSFtitleBarHeight},
  lefttitle=0pt,
  righttitle=0pt,
  toptitle=\ZSFspaceXS,
  bottomtitle=\ZSFspaceXS,
  before title={},
  after title={},
}}

% Rechtsbuendiges Meta-/Kategorie-Tag im Titelbalken einer Titelbox.
% Im Titel-Argument verwenden, z.B.:
%   \begin{defbox}[Satz von Fubini\ZSFTitleTag{Integration}] ... \end{defbox}
\newcommand{\ZSFTitleTag}[1]{\hfill{\ZSFfontTitleTag #1}}

% Gemeinsamer Basis-Stil fuer alle Panel-Boxen mit Titelbalken
\tcbset{zsfpanelbox/.style={
  enhanced jigsaw,
  colback=white,
  colbacktitle=\BoxTitleColor,
  colframe=\BoxTitleColor,
  coltitle=\BoxTitleTextColor,
  boxrule=0.5pt,
  titlerule=0pt,
  boxsep=0pt,
  arc=1pt, outer arc=1pt,
  before skip=\ZSFBoxBeforeSkip,
  after skip=\ZSFspaceS,
  before={\ZSFConsumeAfterSubsectionBar},
  zsftitlebar,
}}

% Gemeinsamer Basis-Stil fuer alle Titelboxen (defbox, tablebox, figbox)
\tcbset{zsftitlebox/.style={
  enhanced jigsaw,
  breakable=false,
  colback=\BoxBackColor,
  colbacktitle=\BoxTitleColor,
  colframe=\BoxFrameColor,
  coltitle=\BoxTitleTextColor,
  boxrule=0.5pt,
  titlerule=0pt,
  boxsep=0pt,
  before skip=\ZSFBoxBeforeSkip,
  after skip=\ZSFspaceS,
  before={\ZSFConsumeAfterSubsectionBar},
  zsftitlebar,
  before upper={\parskip=\ZSFspaceS\relax},
}}

\newtcolorbox{defbox}[1][]{zsftitlebox, before skip=\ZSFBoxBeforeSkip, left=\ZSFspaceS, right=\ZSFspaceS, top=\ZSFspaceS, bottom=\ZSFspaceS, title={#1}}
\newtcolorbox{tablebox}[1][]{zsftitlebox, unbreakable, title={#1}}
\newtcolorbox{figbox}[1][]{zsftitlebox, unbreakable, title={#1}}

% Randlose, modulare Tabelle (für zentrierte Formeln, nützt volle Box aus)
\newtcolorbox{fulltablebox}[2][]{
  zsfpanelbox,
  colframe=black!20, % Hellerer Rahmen für die äußere Tabellenbox
  breakable=false, % Tabellen lieber nicht umbrechen
  title={#1},
  boxsep=0pt, left=0pt, right=0pt, top=\ZSFtextMinPad, bottom=\ZSFtextMinPad,
  tabularx*={\arrayrulecolor{\ZSFtableRuleColor}\setlength{\arrayrulewidth}{\ZSFtableRuleWidth}}{#2},
  zsftitlebar,      % Restore title spacing broken by tabularx defaults
  boxrule=0.5pt     % Restore boxrule broken by tabularx defaults
}

% Falls kein Titel, aber volle Breite & zentriert gesucht wird (Randlose Box)
\newtcolorbox{fulllessbox}[1]{
  enhanced jigsaw,
  breakable=false,
  colback=white,
  colframe=white, % Unsichtbarer Rand
  boxrule=0pt,
  arc=0pt, outer arc=0pt,
  before skip=\ZSFspaceS,
  after skip=\ZSFspaceS,
  boxsep=0pt, left=0pt, right=0pt, top=\ZSFtextMinPad, bottom=\ZSFtextMinPad,
  tabularx*={\arrayrulecolor{\ZSFtableRuleColor}\setlength{\arrayrulewidth}{\ZSFtableRuleWidth}}{#1},
  boxrule=0pt       % Restore boxrule broken by tabularx defaults
}

\newenvironment{longformula}{\[\begin{aligned}}{\end{aligned}\]}
\newtcolorbox{formulabox}[1][]{
  enhanced jigsaw,
  colback=\FormulaboxBackColor,
  colframe=black,  % 1px Rand schwarz
  boxrule=0.5pt,   % Ca. 1px Strichdicke in LaTeX
  left=\ZSFspaceS,
  right=\ZSFspaceS,
  top=\ZSFspaceS,
  bottom=\ZSFspaceS,
  before skip=\ZSFBoxBeforeSkip,
  after skip=\ZSFspaceS,
  before={\ZSFConsumeAfterSubsectionBar},
  before upper={\parskip=\ZSFspaceS\centering},
  after upper={\par},
  middle=\ZSFspaceSep,
  #1
}

\newtcolorbox{goalbox}[1][]{
  zsfpanelbox,
  breakable=false,
  colback=white,
  left=\ZSFspaceS,
  right=\ZSFspaceS,
  top=\ZSFtextMinPad,
  bottom=\ZSFtextMinPad,
  title={#1},
  before upper={\parskip=\ZSFspaceS\centering},
  after upper={\par},
}

\newtcbox{\GoalConditionBox}{%
  on line,
  enhanced,
  boxrule=1.5pt,
  colback=white,
  colframe=\chaptercolorlight!22,
  borderline={0.2pt}{0pt}{black},
  arc=1pt,
  left=\ZSFspaceXS,
  right=\ZSFspaceXS,
  top=\ZSFspaceXS,
  bottom=\ZSFspaceXS,
  nobeforeafter,
  tcbox raise base,
}

\newtcbox{\GoalTargetBox}{%
  on line,
  enhanced,
  boxrule=0.5pt,
  colback=\chaptercolorlight!22,
  colframe=\chaptercolor,
  arc=1pt,
  left=\ZSFspaceXS,
  right=\ZSFspaceXS,
  top=\ZSFspaceXS,
  bottom=\ZSFspaceXS,
  nobeforeafter,
  tcbox raise base,
}

\newcommand{\GoalCondition}[1]{\par\centering\GoalConditionBox{\ZSFfontNote #1}\par}
\newcommand{\GoalStep}[1]{\par\centering$\displaystyle #1$\par}
\newcommand{\GoalTarget}[1]{\par\centering\GoalTargetBox{$\displaystyle #1$}\par}
\newcommand{\ZSFDerivationCase}[3]{%
  \GoalCondition{#1}%
  \GoalStep{#2 =}%
  \GoalTarget{#3}%
}

\newcommand{\ZSFDerivationInlineCase}[3]{%
  \GoalCondition{#1}%
  \GoalStep{#2 = \GoalTargetBox{$\displaystyle #3$}}%
}

% Kompaktere Variante mit reduziertem vertikalen Abstand zwischen Bedingung und Umformung
\newcommand{\ZSFDerivationInlineCaseCompact}[3]{%
  {%
    \setlength{\parskip}{0pt}%
    \par\raggedright%
    \tcbox[
      on line,
      enhanced,
      boxrule=1.5pt,
      colback=white,
      colframe=\chaptercolorlight!22,
      borderline={0.2pt}{0pt}{black},
      arc=1pt,
      boxsep=0pt,
      left=0pt,
      right=0pt,
      top=0pt,
      bottom=0pt,
      nobeforeafter,
      tcbox raise base
    ]{\ZSFfontNote #1}\par
    \vspace{0pt}%
    \centering$\displaystyle #2 = \GoalTargetBox{$\displaystyle #3$}$\par
    \vspace{-0.5pt}%
  }%
}

\newcommand{\GoalCard}[4][]{%
  \begin{goalbox}[#1]
    \GoalCondition{#2}
    #3
    \GoalTarget{#4}
  \end{goalbox}%
}

\newcommand{\GoalPair}[8]{%
  \begin{tcbraster}[raster columns=2, raster column skip=\ZSFspaceM, raster row skip=0pt, raster equal height]
    \GoalCard[#1]{#2}{#3}{#4}
    \GoalCard[#5]{#6}{#7}{#8}
  \end{tcbraster}%
}

% \formulanote: Note-Text unter Formeln - mit parskip-Kompensation
\newcommand{\formulanote}[1]{%
  \par\vspace{-\parskip}\vspace{\ZSFspaceS}%
  \noindent{{\ZSFfontNote #1}}%
  \par\vspace{-\parskip}\vspace{\ZSFspaceS}%
}

% Dünne, graue Trennlinie für Formel-Boxen mit mehreren Elementen
\newcommand{\formulasep}{%
  \par\vspace{-\parskip}\vspace{\ZSFspaceSep}%
  \noindent\textcolor{black!30}{\rule{\linewidth}{0.5pt}}%
  \par\vspace{-\parskip}\vspace{\ZSFspaceSep}%
}


% \runintext: Inline-Text für Tabellen etc. - mit parskip-Kompensation
\newenvironment{runintext}{\par\vspace{-\parskip}\vspace{\ZSFspaceS}\noindent\ignorespaces}{\par}




% ------------------------------------------------------------
% Aussagen (statementbox) und Verfahren (procedure)
% ------------------------------------------------------------
% statementbox: dezente Box mit linkem Farbbalken in Kapitelfarbe.
% Optional fetter Titel in Kapitelfarbe. Fuer Saetze, Definitionen
% und sonstige in sich geschlossene Aussagen, wenn `defbox` mit
% Titelleiste zu schwer wirkt.
\tcbset{zsfstatementstyle/.style={
  enhanced jigsaw,
  breakable=false,
  colback=white,
  colframe=white,
  frame hidden,
  borderline west={\ZSFstatementBarWidth}{0pt}{\StatementBarColor},
  boxrule=0pt,
  arc=0pt, outer arc=0pt,
  left=\ZSFstatementLeftPad,
  right=\ZSFspaceS,
  top=\ZSFspaceXS,
  bottom=\ZSFspaceXS,
  before skip=\ZSFBoxBeforeSkip,
  after skip=\ZSFspaceS,
  before={\ZSFConsumeAfterSubsectionBar},
}}

\NewDocumentEnvironment{statementbox}{O{}}
  {\begin{tcolorbox}[zsfstatementstyle]%
   \if\relax\detokenize{#1}\relax\else
     {\bfseries\color{\StatementTitleColor}#1}\par\vspace{\ZSFspaceXS}%
   \fi\ignorespaces}
  {\end{tcolorbox}}

% procedure: nummerierte Schritt-Liste in einer statementbox.
% Jeder Schritt = fetter Titel + eigene Zeile mit Beschreibung.
% Verwendung: \ProcStep{Titel}{Beschreibung}
\NewDocumentEnvironment{procedure}{O{}}
  {\begin{statementbox}[#1]%
   \begin{enumerate}[leftmargin=1.6em,labelsep=0.4em,itemsep=\ZSFprocItemSep,topsep=0pt,parsep=0pt,partopsep=0pt]}
  {\end{enumerate}\end{statementbox}}

\newcommand{\ProcStep}[2]{\item {\bfseries #1}\\[\ZSFprocStepGap]#2}

% factlist: rahmenlose Liste von Definitions-/Eigenschaftsfakten.
% Farbiger Bullet, fetter Lead-in, em-dash, Beschreibung.
% Verwendung: \ZSFFact{Lead-in}{Beschreibung} - oder rohes \item.
\NewDocumentEnvironment{factlist}{}
  {\begin{itemize}[
     leftmargin=1.2em,
     labelsep=0.4em,
     label={\textcolor{\StatementBarColor}{\textbullet}},
     itemsep=\ZSFprocItemSep,
     topsep=\ZSFspaceXS,
     parsep=0pt,
     partopsep=0pt
   ]}
  {\end{itemize}}

% propertylist: gruppierte Charakterisierungen in einer statementbox.
% Wie factlist, aber zusammengefasst unter optionalem Titel.
\NewDocumentEnvironment{propertylist}{O{}}
  {\begin{statementbox}[#1]\begin{itemize}[
     leftmargin=1.2em,
     labelsep=0.4em,
     itemsep=\ZSFprocItemSep,
     topsep=0pt,
     parsep=0pt,
     partopsep=0pt
   ]}
  {\end{itemize}\end{statementbox}}

\newcommand{\ZSFFact}[2]{\item \textbf{#1}\,---\,#2}

\newtcbox{\ZSFdanger}{colback=red!80!black, colframe=red!80!black, coltext=white, on line, boxsep=\ZSFspaceXS, left=\ZSFspaceS, right=\ZSFspaceS, top=\ZSFspaceXS, bottom=\ZSFspaceXS, boxrule=0pt, arc=1pt, fontupper=\sffamily\bfseries}

% --- Nebenläufiges Bild/Text-Layout (für "Ebene Kurven", Tabellen mit Bildern etc.) ---
\newtcolorbox{splitbox}[1][0.3]{
  blankest,
  sidebyside,
  sidebyside align=top,
  sidebyside gap=\ZSFspaceM,
  lefthand width=#1\linewidth,
  lower separated=false,
  before skip=\ZSFspaceS,
  after skip=\ZSFspaceS,
  boxsep=0pt, left=0pt, right=0pt, top=0pt, bottom=0pt
}

% --- Zentriertes Item-Heading für Listen mit Trennstrich (für Kurven, Profile etc.) ---
\newcommand{\ZSFItemHeading}[1]{%
  \par\vspace{-\parskip}\vspace{\ZSFspaceM}%
  {\centering\textbf{#1}\par}\vspace{-\parskip}\vspace{\ZSFspaceS}%
  \noindent\rule{\linewidth}{0.5pt}\par\vspace{-\parskip}\vspace{\ZSFspaceS}%
}



% Kompakte Grid-Box als tabularray-Container
\newtcolorbox{compactgridbox}[1][]{
  zsfpanelbox,
  colframe=black!20,
  breakable=false,
  title={#1},
  boxsep=0pt, left=0pt, right=0pt, top=\ZSFtextMinPad, bottom=\ZSFtextMinPad,
  zsftitlebar,
  boxrule=0.5pt
}

% Gemeinsamer tabularray-Stil fuer alle Valuegrids
\newcommand{\ZSFvaluegridCommon}{
      colsep=4pt,
      rowsep=3pt,
      hlines,
      vlines,
      rows={valign=m}
}

% Parametrisiertes Valuegrid: #1 = Anzahl Datenspalten, #2 = optionaler Titel
% Erzeugt automatisch `l *{#1}{c}` als colspec
\NewDocumentEnvironment{valuegrid}{m O{}}
  {\begin{compactgridbox}[#2]\begin{tblr}{colspec={l *{#1}{c}},\ZSFvaluegridCommon}}
  {\end{tblr}\end{compactgridbox}}

% Freie colspec-Variante fuer Sonderfaelle
\NewDocumentEnvironment{valuegridcustom}{m O{}}
  {\begin{compactgridbox}[#2]\begin{tblr}{colspec={#1},\ZSFvaluegridCommon}}
  {\end{tblr}\end{compactgridbox}}

% --- Legacy-Aliase (Abwaertskompatibilitaet) ---
\NewDocumentEnvironment{valuegridtwo}{O{}}{\begin{valuegrid}{2}[#1]}{\end{valuegrid}}
\NewDocumentEnvironment{valuegridthree}{O{}}{\begin{valuegrid}{3}[#1]}{\end{valuegrid}}
\NewDocumentEnvironment{valuegridfour}{O{}}{\begin{valuegrid}{4}[#1]}{\end{valuegrid}}
\NewDocumentEnvironment{valuegridfive}{O{}}{\begin{valuegrid}{5}[#1]}{\end{valuegrid}}
\NewDocumentEnvironment{valuegridsix}{O{}}{\begin{valuegrid}{6}[#1]}{\end{valuegrid}}
\NewDocumentEnvironment{valuegridseven}{O{}}{\begin{valuegrid}{7}[#1]}{\end{valuegrid}}

